%! Boxplots and Comparing Distributions — Unit 1, Lesson 1.6 — Student Packet Cover Sheet
% GRADUAL-RELEASE cover: four packet rows (Warm-Up · Guided Notes & Practice · AP Practice · Homework).
% The EFFL spoiler rule is GONE — the learning targets name the formal vocabulary and bold it,
% because the guided notes teach it directly today. The AP Learning Objectives (UNC-1.L through
% UNC-1.O) stay in the teacher-facing plan.
% NAMESTRIP: \namedateperiod appears HERE and nowhere else in the lesson.
% The `Keep in Mind' box carries the lesson's IDEAS, never its process — no narration of the
% gradual release.
\documentclass[12pt]{article}
\usepackage{apstats-article}
\usepackage{apstats-boxes}
\usepackage{ltablex}
\keepXColumns

\begin{document}

% Full-bleed navy banner. \coverbanner MEASURES the title block and sizes the
% band to it, so a lesson title long enough to wrap grows the band rather than
% running out the bottom of it.
\coverbanner{Unit 1}{Lesson 1.6 \quad Boxplots and Comparing Distributions}

\namedateperiod
\vspace{0.18in}

\boxguard[14]
\begin{learningtargetbox}
\small
\begin{enumerate}[leftmargin=1.5em, itemsep=5pt]
  \item \ldots{} report the \textbf{five-number summary} of a data set, and point to each of
        those five numbers on a \textbf{boxplot}.
  \item \ldots{} explain why every section of a boxplot holds the \textbf{same quarter} of the
        data, however wide it looks --- and say what the width of a section really reports.
  \item \ldots{} use the $1.5 \times \text{IQR}$ rule to decide whether a value is an
        \textbf{outlier}, and say how a \textbf{modified boxplot} draws it and where the
        \textbf{whisker} stops.
  \item \ldots{} compare two distributions from \textbf{parallel (side-by-side) boxplots} on
        center, variability, and unusual features --- in context, with comparative language.
\end{enumerate}
\end{learningtargetbox}

\vspace{0.14in}

\boxguard[14]
\begin{tocbox}
\small
\renewcommand{\arraystretch}{1.5}
\begin{tabularx}{\linewidth}{@{} c l X r @{}}
\rowcolor{sky}
\textbf{\#} & \textbf{Component} & \textbf{Description} & \textbf{Score} \\
\midrule
1 & Warm-Up & Five numbers off an ordered list, a counting drill, and one $1.5 \times
                 \text{IQR}$ check & \blank{1.2cm} \\
2 & Guided Notes \& Practice & Five numbers and the picture drawn from them, four sections that
                 all hold a quarter, and two bus routes on one axis --- then \emph{Two Science
                 Sections} worked together & \blank{1.2cm} \\
% AP Practice is assigned but NOT scored: its score cell prints NA, never a \blank{}.
% Row 4 is the lesson's GRADED practice. Paper homework is the DEFAULT for this course; on the
% occasions DeltaMath is assigned instead, that replaces row 4 and its score cell is struck.
3 & AP Practice & Weekday and weekend tips: AP-style multiple choice and one free-response set
                 --- practice, not a grade & \textbf{NA} \\
4 & Homework & Five numbers matched to a plot, a claim to correct, one outlier check, a
                 two-shop comparison, and one multiple-choice item --- \textbf{scored, due the
                 first class after two study halls} & \blank{1.2cm} \\
\midrule
  & \multicolumn{2}{r}{\textbf{Total}} & \blank{1.2cm} \\
\end{tabularx}
\end{tocbox}

\vspace{0.14in}

\boxguard[8]
\begin{remindbox}
\small
A boxplot is five numbers on a number line and nothing more. Its four sections look wildly
different in width, and \textbf{every one of them holds exactly a quarter of the data} --- so
width is telling you how far apart those values are, never how many there are. \textbf{Count
before you conclude.} And remember what those five numbers cost you: a gap, a clump, or a second
peak disappears inside the box, so when somebody asks a question about shape, ask for a dotplot.
When you compare two groups, a description of one followed by a description of the other is not a
comparison --- the comparing word is the part that earns the credit.
\end{remindbox}

\end{document}
